\section*{Chapitre \ref{ch:ml:reinforcement-learning-foundations} --- Fondements de l'apprentissage par renforcement}

\begin{solution}{exo:ml:reinforcement-learning-foundations:1}
La seule action consiste à vendre les 3 lots~: coût $\eta\times3^2$, et rien n'est détenu ensuite. Liquidité mince~: $-2\times9 = -18$ points de
base~; liquidité profonde~: $-0.5\times9 = -4.5$.
\end{solution}

\begin{solution}{exo:ml:reinforcement-learning-foundations:2}
\omterm{def:ml:reinforcement-learning-foundations:td}{Q-learning}~: $-10 + 0.1(-3 - 5 + 10) = -9.8$. SARSA~: $-10 + 0.1(-3 - 8 + 10) = -10.1$~: elle évalue l'action selon ce que la
\omterm{def:ml:reinforcement-learning-foundations:mdp}{politique} d'exploration va réellement faire ensuite.
\end{solution}

\begin{solution}{exo:ml:reinforcement-learning-foundations:3}
Les coûts de détention valent $0.2q^2$ par période après chaque vente~: avec un lot par période, $0.2(9^2 + 8^2 + \dots + 1^2) = 57$ points de base de
risque sur 10 lots, contre un impact d'au plus 20. Le desk avance les ventes (3 lots d'abord) parce que le risque
l'emporte ici sur l'impact~; l'optimum avance davantage les ventes lorsque la liquidité est profonde.
\end{solution}

\begin{solution}{exo:ml:reinforcement-learning-foundations:4}
SARSA apprend la valeur de la \omterm{def:ml:reinforcement-learning-foundations:mdp}{politique} qu'il suit, exploration comprise, si bien que sa \omterm{def:ml:reinforcement-learning-foundations:mdp}{politique} gloutonne évite les actions coûteuses
lorsque l'exploration tourne mal~; au début, cette prudence est utile. \omterm{def:ml:reinforcement-learning-foundations:td}{Q-learning} apprend directement les valeurs optimales et, une
fois celles-ci apprises, sa \omterm{def:ml:reinforcement-learning-foundations:mdp}{politique} gloutonne est la \omterm{def:ml:reinforcement-learning-foundations:mdp}{politique} optimale~; celle de SARSA reste optimale pour une \omterm{def:ml:reinforcement-learning-foundations:mdp}{politique} exploratoire, un peu plus prudente que
nécessaire.
\end{solution}

\begin{solution}{exo:ml:reinforcement-learning-foundations:5}
Le poids d'un épisode est un produit de dix rapports. Avec $k$ actions possibles, une action journalisée que la cible préfère et
que le desk ne préfère pas a pour rapport $(0.9 + 0.1/k)/(0.3/k)\approx3k$~; une action que le desk préfère et que la cible ne préfère pas a pour rapport
$(0.1/k)/(0.7 + 0.3/k)$, bien inférieur à un. Sur 2\,000 épisodes journalisés, la
médiane du poids vaut $2.8\times10^{-5}$, le 99\up{e}~centile vaut 1.56 et le plus grand 797~: un seul épisode porte 78~\% du
poids total. L'estimation revient en pratique à une moyenne sur une poignée d'épisodes.
\end{solution}

\begin{solution}{exo:ml:reinforcement-learning-foundations:6}
La statistique $t$ mesure le bruit du simulateur, et non son erreur~: on peut prolonger une simulation d'une année autant
que l'on souhaite. La question est de savoir si le simulateur est juste sur les points que l'agent exploite. Évaluer l'agent sur des journaux réels
(\omterm{def:ml:reinforcement-learning-foundations:ope}{évaluation hors politique}), dans des simulateurs aux hypothèses différentes et lors d'un essai réel limité face au dispositif en place.
\end{solution}

\begin{solution}{exo:ml:reinforcement-learning-foundations:7}
Avec la ponction de liquidité, le coût optimal passe de 2.59 à 3.01 points de base par lot, et la première vente en liquidité mince passe
de 2 à 3 lots (5 en liquidité profonde, inchangé). Vendre dès maintenant rend les périodes suivantes plus susceptibles d'être minces, si bien qu'attendre vaut
moins~; l'optimum vend davantage tant qu'il le peut.
\end{solution}

\begin{solution}{exo:ml:reinforcement-learning-foundations:8}
Pour un pas~: $\hat V_{\mathrm{DR}} = \hat V(x) + w(g - \hat Q(x, u))$ avec $w = \pi_e(u\mid x)/\pi_b(u\mid x)$. En prenant
l'espérance sous $\pi_b$ et sous la loi de la \omterm{def:ml:reinforcement-learning-foundations:mdp}{récompense}, $\E[w(g - \hat Q(x, u))] = \sum_u\pi_e(u\mid x)(Q^{\pi_e}(x, u) -
\hat Q(x, u))$, qui s'annule lorsque $\hat Q = Q^{\pi_e}$, ce qui laisse $\hat V(x) = V^{\pi_e}(x)$. Si en revanche les poids sont
justes et $\hat Q$ est faux, la même expression vaut $V^{\pi_e}(x) - \hat V(x)$ et la somme vaut à nouveau $V^{\pi_e}(x)$~:
sans biais si l'une des deux parties est juste.
\end{solution}

\begin{solution}{pb:ml:reinforcement-learning-foundations:1}
\textbf{Partie I.}
\begin{enumerate}
\item États (période, lots restants, liquidité)~; actions~: lots vendus~; \omterm{def:ml:reinforcement-learning-foundations:mdp}{récompense}~: l'opposé de la somme de l'impact $\eta u^2$ et du risque $0.2q^2$~;
la liquidité persiste avec une probabilité de 0.8~; tout ce qui reste est vendu à la fin.
\item 2.59 points de base par lot~; 5 lots d'abord en liquidité profonde, 2 en liquidité mince.
\item 3.27 et 6.76.
\item L'état est minuscule et les probabilités de transition sont connues~; un marché a des prix continus, une dynamique inconnue et
des états qu'il ne révèle jamais.
\end{enumerate}
\textbf{Partie II.}
\begin{enumerate}[resume]
\item Pour 1\,000 et 10\,000 épisodes~: \omterm{def:ml:reinforcement-learning-foundations:td}{Q-learning} 1.41 et 0.044, SARSA 0.28 et 0.051, REINFORCE 0.31 et 0.067, \omterm{def:ml:reinforcement-learning-foundations:pg}{méthode
acteur--critique} 0.11 et 0.009 points de base par lot.
\item De grands pas poussent le softmax vers une action avant que sa valeur ne soit bien estimée~; la \omterm{def:ml:reinforcement-learning-foundations:mdp}{politique} cesse alors d'échantillonner les
autres possibilités et ne peut pas apprendre qu'elles sont meilleures.
\item Il se met à jour à chaque pas à partir de l'erreur de différence temporelle, avec une variance plus faible que celle du rendement de l'épisode entier.
\item Avec vingt ordres parents par jour, sur 500 jours de négociation~: deux années d'ordres d'un desk pour un problème à 242 états.
\end{enumerate}
\textbf{Partie III.}
\begin{enumerate}[resume]
\item $\varepsilon$-glouton 0.1~: 1\,727~; 0.01~: 2\,026~; UCB avec $c = \sqrt2$~: 1\,325~; avec $c = 1$~: 944 points de base. UCB avec
$c = 1$ est le meilleur.
\item La règle explore si peu que, dans certaines séries, elle se fixe sur un bras qui paraît bon sans l'être, et revisite rarement les autres.
\item Biais et écart-type~: échantillonnage préférentiel 0.20 et 1.75~; échantillonnage préférentiel pondéré 0.49 et 0.63~; modèle seul 0.30 et 0~; doublement
robuste $-0.03$ et 0.29.
\item L'état, l'action et la probabilité avec laquelle la \omterm{def:ml:reinforcement-learning-foundations:mdp}{politique} de comportement l'a choisie, pour chaque décision.
\end{enumerate}
\textbf{Partie IV.}
\begin{enumerate}[resume]
\item \emph{Apprendre à vendre.} Après 10\,000 épisodes, la \omterm{def:ml:reinforcement-learning-foundations:mdp}{politique} gloutonne de \omterm{def:ml:reinforcement-learning-foundations:td}{Q-learning} est à 0.044 point de base par lot
de l'optimum (2.59). Estimations hors \omterm{def:ml:reinforcement-learning-foundations:mdp}{politique} d'une \omterm{def:ml:reinforcement-learning-foundations:mdp}{politique} coûtant 2.77~: échantillonnage préférentiel biaisé de 0.20 avec un écart-type de 1.75, échantillonnage préférentiel pondéré
0.49 et 0.63, doublement robuste $-0.03$ et 0.29.
\item Pour attendre la fin d'une liquidité mince, que le simulateur faisait toujours cesser dès la période suivante~; sur le marché, il a coûté 3.83 contre
3.27 pour le desk, après 2.77 contre 3.66 dans le simulateur.
\item En évaluant dans un simulateur à liquidité persistante, par \omterm{def:ml:reinforcement-learning-foundations:ope}{évaluation hors politique} sur des journaux réels, ou en confrontant
aux données le processus de liquidité du simulateur.
\item Lorsque les actions ne modifient pas les états futurs (router un ordre entre plusieurs algorithmes d'exécution), et que la rétroaction est immédiate.
\item Elle n'utilise aucune hypothèse sur la liquidité, si bien qu'une erreur sur cette hypothèse ne peut pas lui nuire.
\item Les tables deviennent des approximateurs de fonction, l'exploration devient coûteuse et l'évaluation exacte disparaît~; l'évaluation doit
reposer sur des journaux et sur plusieurs simulateurs.
\item Un banc d'essai pour les mécanismes et pour les bogues, et un lieu où comparer des \omterm{def:ml:reinforcement-learning-foundations:mdp}{politiques} sous hypothèses explicites, jamais le
juge final.
\item L'environnement dans lequel elle a été entraînée, erreurs comprises.
\end{enumerate}
\end{solution}

\begin{solution}{iq:ml:reinforcement-learning-foundations:1}
$Q^*(x, u) = \E[g_{t+1} + \gamma\max_{u'}Q^*(x_{t+1}, u')\mid x_t = x, u_t = u]$~: la \omterm{def:ml:reinforcement-learning-foundations:mdp}{récompense} immédiate, plus la valeur actualisée
d'un choix optimal à partir de l'état suivant, moyennée sur les transitions.
\iqlookfor{l'équation et ses termes.}
\end{solution}

\begin{solution}{iq:ml:reinforcement-learning-foundations:2}
La cible du \omterm{def:ml:reinforcement-learning-foundations:td}{Q-learning} utilise la meilleure action suivante (hors \omterm{def:ml:reinforcement-learning-foundations:mdp}{politique}, il apprend les valeurs optimales)~; celle de SARSA utilise l'action suivante effectivement prise
(sur \omterm{def:ml:reinforcement-learning-foundations:mdp}{politique}, il apprend les valeurs de la \omterm{def:ml:reinforcement-learning-foundations:mdp}{politique} d'exploration). Cela compte lorsque l'exploration est coûteuse~: SARSA apprend à éviter les états où
une action exploratoire est dangereuse.
\iqlookfor{cibles hors \omterm{def:ml:reinforcement-learning-foundations:mdp}{politique} et sur \omterm{def:ml:reinforcement-learning-foundations:mdp}{politique}, et un cas où elles diffèrent.}
\end{solution}

\begin{solution}{iq:ml:reinforcement-learning-foundations:3}
\omterm{def:ml:reinforcement-learning-foundations:ope}{Évaluation hors politique} à partir des journaux, avec les probabilités de la \omterm{def:ml:reinforcement-learning-foundations:mdp}{politique} de comportement~: échantillonnage préférentiel, pondéré ou doublement robuste avec
un modèle d'exécutions et d'impact~; vérifier le support (les actions de la nouvelle \omterm{def:ml:reinforcement-learning-foundations:mdp}{politique} doivent avoir été journalisées), indiquer la dispersion et confirmer
par un petit test en conditions réelles.
\iqlookfor{pondération préférentielle, support et correction fondée sur un modèle.}
\end{solution}

\begin{solution}{iq:ml:reinforcement-learning-foundations:4}
Un \omterm{def:ml:reinforcement-learning-foundations:bandit}{bandit contextuel}~: le contexte est l'ordre et l'état du marché, la \omterm{def:ml:reinforcement-learning-foundations:mdp}{récompense} le coût par rapport à un indice de référence, la règle UCB ou
échantillonnage de Thompson avec un modèle de \omterm{def:ml:reinforcement-learning-foundations:mdp}{récompense} par courtier, avec des allocations minimales pour l'apprentissage et des probabilités journalisées pour
l'évaluation ultérieure.
\iqlookfor{formulation en bandit, règle d'exploration et journalisation.}
\end{solution}

\begin{solution}{iq:ml:reinforcement-learning-foundations:5}
$\nabla_\theta J = \E_\pi[\sum_t\nabla_\theta\log\pi_\theta(u_t\mid x_t)(G_t - b(x_t))]$~; soustraire une baseline qui ne
dépend pas de l'action laisse l'espérance inchangée, car $\E[\nabla\log\pi] = 0$, et réduit la variance lorsque $b$
est proche de la valeur.
\iqlookfor{le théorème et la raison pour laquelle la baseline est sans biais.}
\end{solution}

\begin{solution}{iq:ml:reinforcement-learning-foundations:6}
Erreurs du simulateur qu'il exploite (exécutions, position dans la file, décroissance de l'impact, persistance de la liquidité, réactions des autres agents), latence,
non-stationnarité et états rares. Tester dans des simulateurs qui font varier chaque hypothèse, évaluer sur des journaux réels hors \omterm{def:ml:reinforcement-learning-foundations:mdp}{politique}, plafonner ses
actions, et le lancer en petit face au dispositif en place avec un kill switch.
\iqlookfor{exploitation du simulateur et validation par étapes.}
\end{solution}
